\section{原型系统与评估方法}\label{sec:methodology}

\emph{单发送机与双接收链测试台架：}单台 Zynq-7020/AD9363 SDR 发送机以标称 \SI{0}{dBm} 及 \SI{10}{dB} 板端衰减发射，采用有限超时主动轮询直接内存访问（DMA）完成状态作为提交代理，而非校准后的射频（RF）空口时长或能耗；全部 36 轮确认性实验均冻结同一套 FreeRTOS \proto{} 发送固件。两块独立 HackRF One 接收机以 \SI{4}{Msps} 采样（关闭射频功放、固定天线绑定，低噪声放大器（LNA）\SI{16}{dB}；可变增益放大器（VGA）在有线 \texttt{R0} 轮次为 \SI{10}{dB}，无线 \texttt{R1} 轮次为 \SI{16}{dB}），分别送入软件 \ble{} 解码链（监听 \ble{} ch3 或 ch39）与标准偏移正交相移键控（O-QPSK）软件解码链（监听 \stdfour{} ch12）。各链先独立校验其外层帧结构、\proto{} 头及记录 CRC；对于分片 \modelc{} 记录，\ble{} 链须先在同一次尝试内集齐全部 \fragproto{} 分片。确定性净荷字节仅用于解码完整性核验，严禁用于重构对侧漏收字节；已校验记录仅通过 $(\texttt{run\_id},\texttt{command\_id})$ 实施跨链匹配。

% ============================================================
%  表 2（中文同步版）：与 ../tables/tab-p4-main.tex 数值 100% 同步
%  数据源：python/experiments/p5_freeze_20261124/{p5_p4_cells.csv, p5_p4_paired_blocks.csv}
% ============================================================
\begin{table*}[t]
\centering
\caption{36 轮交替区组无线确认性实验的计划尝试内联合命令解析率与单位波形代理成本（\num{3600} 项计划任务，部分缺失条件 \texttt{R1}）。}
\label{tab:p4_main}
\setlength{\tabcolsep}{2.4pt}
\footnotesize
\begin{tabular}{@{}llccccccc@{}}
\toprule
\textbf{条件} & \textbf{模式} & $K$ & \textbf{首次} & \textbf{三区组} & \textbf{合并} & \textbf{均值$\pm$SD} & \textbf{单位成本} & \textbf{配对差 $\modejc{}-\text{模式}$ (pp)} \\
 & & & ($/300$) & B1/B2/B3 & ($/300$) & (\%) & (\si{\micro\second}/任务) & B1/B2/B3 (均值; 95\% $t$-CI) \\
\midrule
\multirow{3}{*}{A $(0,0)$, $K{=}1$}
  & \modejc{}              & 1 & 205/300 & 63 / 78 / 64 & 205/300 & $68.3 \pm 8.4$  & \num{2411.7}  & --- \\
  & \modesc{}              & 1 & 212/300 & 66 / 81 / 65 & 212/300 & $70.7 \pm 9.0$  & \num{2332.1}  & $-3/{-}3/{-}1$（$-2.33$；$[-5.20, +0.54]$） \\
  & \modelc{}$^{\mathrm{a}}$ & 1 & 204/300 & 51 / 90 / 63 & 204/300 & $68.0 \pm 20.0$ & \num{2141.2}  & $+12/{-}12/{+}1$（$+0.33$；$[-29.51, +30.18]$） \\
\midrule
\multirow{3}{*}{A $(0,0)$, $Q{=}6000$}
  & \modejc{}              & 3 & 224/300 & 95 / 100 / 99  & 294/300 & $98.0 \pm 2.6$ & \num{5044.9}  & --- \\
  & \modesc{}              & 3 & 227/300 & 99 / 100 / 100 & 299/300 & $99.7 \pm 0.6$ & \num{4960.5}  & $-4/{+}0/{-}1$（$-1.67$；$[-6.84, +3.50]$） \\
  & \modelc{}$^{\mathrm{a}}$ & 4 & 201/300 & 99 / 100 / 96  & 295/300 & $98.3 \pm 2.1$ & \num{5922.7}  & $-4/{+}0/{+}3$（$-0.33$；$[-9.06, +8.39]$） \\
\midrule
\multirow{3}{*}{B $(26,0)$, $K{=}1$}
  & \modejc{}              & 1 & 238/300 & 72 / 86 / 80 & 238/300 & $79.3 \pm 7.0$  & \num{2339.5}  & --- \\
  & \modesc{}              & 1 & 216/300 & 75 / 78 / 63 & 216/300 & $72.0 \pm 7.9$  & \num{2577.8}  & $-3/{+}8/{+}17$（$+7.33$；$[-17.55, +32.22]$） \\
  & \modelc{}$^{\mathrm{a}}$ & 1 & 86/300  & 43 / 36 / 7  & 86/300  & $28.7 \pm 19.1$ & \num{10130.2} & $+29/{+}50/{+}73$（$+50.67$；$[-4.00, +105.34]^{\mathrm{b}}$） \\
\midrule
\multirow{3}{*}{B $(26,0)$, $Q{=}6000$}
  & \modejc{}              & 3 & 195/300 & 99 / 100 / 92 & 291/300 & $97.0 \pm 4.4$  & \num{5740.2}  & --- \\
  & \modesc{}              & 3 & 224/300 & 99 / 99 / 99  & 297/300 & $99.0 \pm 0.0$  & \num{5624.2}  & $+0/{+}1/{-}7$（$-2.00$；$[-12.83, +8.83]$） \\
  & \modelc{}$^{\mathrm{a}}$ & 2 & 56/300  & 76 / 0 / 17   & 93/300  & $31.0 \pm 39.9$ & \num{18735.5} & $+23/{+}100/{+}75$（$+66.00$；$[-31.58, +163.58]^{\mathrm{b}}$） \\
\bottomrule
\end{tabular}
\par\smallskip
\raggedright\scriptsize
每单元合并 $n{=}3$ 独立交替区组 $\times$ $100$ 计划任务（每单元 $300$ 项；四分类结果见图~\ref{fig:p4_tradeoff}a）。双侧在至多 $K$ 次计划尝试内均校验通过即计为联合解析（不含截止期前按时交付审计）；\modelc{} 须在单次尝试内集齐全部分片。每单元提交/未用波形代理（\si{\micro\second}）：A-$K{=}1$ \num{494400}/$0$（\modejc{},\modesc{}）、\num{436800}/$0$（\modelc{}）；A-$Q$ \num{1483200}/\num{316800}（\modejc{},\modesc{}）、\num{1747200}/\num{52800}（\modelc{}）；B-$K{=}1$ \num{556800}/$0$（\modejc{},\modesc{}）、\num{871200}/$0$（\modelc{}）；B-$Q$ \num{1670400}/\num{129600}（\modejc{},\modesc{}）、\num{1742400}/\num{57600}（\modelc{}；B2 轮次 2743 解析 $0/100$，单轮成本记 \texttt{N/A}）。$^{\mathrm{a}}$\modejc{}/\modesc{}：BLE ch3（\SI{2410}{MHz}）；\modelc{}：ch39（\SI{2480}{MHz}，配置级）。$^{\mathrm{b}}$$\text{df}{=}2$ $t$-CI（$t_{0.975,2}{=}4.303$）跨零并超出 $\pm 100$\,pp。
\end{table*}


\emph{实验矩阵与审计指标：}有线 \texttt{R0} 实验（共 $18$ 轮）核验表~\ref{tab:budget} 全部 10 个分支、同相/正交（IQ）共享偏移解码、等预算桥接、固定次信道时间运行及离线故障注入。在通过 $12$ 轮先导实验预选出、无溢出、无采集截断且无削波的 \SI{0.10}{m} 天线间距部分丢包条件（\texttt{R1}）下，确认性评估共执行 $36$ 轮交替无线实验（$3$ 区组 $\times$ $12$ 条件 $\times$ $100$ 计划任务 $=$ \num{3600} 项任务，表~\ref{tab:p4_main}），在 A 点 $(0,0)$、B 点 $(26,0)$ 与 $\{\modejc{},\modesc{},\modelc{}\}$ 上对比单次尝试（$K{=}1$）以及预先选定的 \SI{6}{ms} 单任务波形预算（$Q{=}\SI{6000}{\micro\second}$，按各模式可容纳的完整尝试次数分配）。全部比率采用预登记分母（保留未成功任务，分母为零记 \texttt{N/A}）并报告计划尝试内联合解析指标：只要两链均在 $K$ 次尝试内校验通过（即便位于不同槽位），该计划任务也记为一次 \texttt{both}（而 \modelc{} 必须在单次尝试内集齐全部 \fragproto{} 分片），其余归入 \texttt{ble\_only}、\texttt{zigbee\_only} 与 \texttt{neither}。单位成本以总提交波形代理（含重发、纯时间更新及残片）除以联合成功任务数；汇总报告区组均值 $\pm$ 标准差（SD），配对比较报告百分点（pp）差值及自由度 $\text{df}{=}2$（$t_{0.975,2}{=}4.303$）的 95\% 置信区间（CI）。
